\section{Finite lower bounds}\label{sec:finite}
A claim $(n,r,L)$ excludes every network of size at most $L$ in the enlarged
model. For a concrete prefix $P$ of length $p$ put
\[
 A=P(H_w),\qquad B=P(H_{w+1}),\qquad k=L-p.
\]
A \emph{completion} of $(A,B,k)$ is an oriented suffix of length at most
$k$ with a fixed output zero on all rows of $A$ and one on all rows of $B$.
By \cref{prop:zero} this is equivalent to completing $P$. The definition
refers to the two image sets alone.

\paragraph{Local exclusions.}
Every comparator has at most two preimages of any specified vector.
Outside the two ports the input is fixed, and inside them there are at
most two orders of the output pair. A layer with more than $2^k$ distinct
images cannot become a singleton in $k$ steps, as \cref{lem:rigid} requires.

For each original position and sign trace the route a strict sentinel would
take in $P$. Let $b$ be a previously proved lower bound for the rank after
deletion, or zero if none is available. On a valid rank deletion, a full
completion of total size at most $L$ must have at most $D=L-b$ encounters
on that path by \cref{lem:delete}. When deletion would remove the selected
rank itself, use only $D=L$, because each encounter uses one gate.
If a trace exceeds $D$, the prefix is impossible. If it attains $D$, its
present wire cannot be touched again. Variation of that wire \emph{within}
either layer therefore excludes the prefix. A difference solely between
layers does not suffice.

\paragraph{Enumeration.}
Try one representative $(i,j)$ with $i<j$ for each unordered pair.
The opposite orientation differs exactly by swapping the two output
labels. Conjugate every subsequent oriented gate and rename the terminal
coordinate by this transposition. This preserves existence and length of
a suffix. If the two discordant port patterns $10$ and $01$ do not both
occur in $A\cup B$, the gate acts there as one fixed permutation
(identity or transposition). Omitting it and conjugating the suffix gives
the same separation in one fewer step. The permutation need only agree on
these rows, since separation there implies correctness everywhere. Thus
a shortest nonempty completion can start with an enumerated pair. If the
representative violates a path cap, its exclusion rules out its reverse
as well, by conjugacy of completed networks.

\paragraph{Memoization.}
Permute both layers together, sorting wire labels stably by the pair of
column sums in $A,B$. Compare the two resulting \emph{entire} images.
Equal keys exhibit a specific common permutation. Ties can cause equivalent
states to have different keys, which only forgoes reuse. If a completed
exclusion with the same key allowed $k'\ge k$, conjugation transfers any
putative suffix to that earlier state. Even if the two prefixes have
different extreme traces, such a suffix appended to the \emph{earlier}
prefix would contradict its established exclusion at total budget
$|P_{\mathrm{earlier}}|+k'\le L$. A partial search of the earlier state
would not justify this inference.

\begin{proposition}[Replay soundness]\label{prop:replay}
Assume the smaller-width bounds used in replay. If replay accepts
$(n,r,L)$, then $C(n,r)>L$.
\end{proposition}
\begin{proof}
Associate to each closed node the proposition that its image pair has no
completion within its remaining budget. Check first that there is no
separating wire already, excluding the empty suffix. Cardinality and
frozen-wire leaves establish the proposition by the local exclusions.
A reference uses only a proposition already closed in the same claim.
For an expansion, choose a shortest hypothetical suffix. Its first gate
has a useful representative by the enumeration argument. A path violation
contradicts deletion; otherwise replay requires and checks the corresponding
child, whose proposition excludes the rest of the suffix. This proves the
parent proposition. Induct in order of node \emph{closure}, inserting an
expanded state only after its last child. The root has the original layers
and budget $L$. Induction over claims establishes all reported bounds:
the table is reset for each claim and cross-claim dependencies have
strictly smaller width.
\end{proof}

\begin{center}
\begin{minipage}{0.96\linewidth}\small
\textbf{Replay of an implicit preorder tree}\\[2pt]
\texttt{visit}$(P,A,B,k)$: consume one opcode; reject a current separator.\\
\hspace*{1em}1: verify $\max(|A|,|B|)>2^k$; return.\\
\hspace*{1em}2: find a closed common-permutation key with budget $\ge k$; return.\\
\hspace*{1em}3: find a saturated trace on a within-layer varying wire; return.\\
\hspace*{1em}0: require $k>0$; enumerate $i<j$ in lexicographic order.\\
\hspace*{2em}Skip fixed permutations and pairs exceeding a trace allowance.\\
\hspace*{2em}For each remaining $g$, call \texttt{visit}$(Pg,gA,gB,k-1)$.\\
\hspace*{2em}Record the closed key and budget.\\
Reject failed predicates, invalid opcodes, missing bytes and trailing nodes.
\end{minipage}
\end{center}

\section{The exact size}
\begin{proof}[Proof of \cref{thm:main}]
Apply \cref{prop:replay} to $(9,4,18)$ and $(9,5,18)$. Both nine-input
ranks require at least 19 comparators. Duality gives $C(9,6)\ge19$.
The layer of weight five has respectively five and six zeros at widths ten
and eleven. It is a boundary layer for each of $(10,5),(10,6),(11,6)$.
Applying \cref{cor:step} to the first two gives a lower bound of 22 in
both cases. Applying it to the third gives 25. This suffices for the
fixed-wire, lower-index-minimum model. For the opposite inequality, the
25 comparators in \cref{fig:net} yield the required bit on wire 5 on all
$2^{11}$ binary inputs by exhaustive evaluation. Apply \cref{prop:zero}.
\end{proof}
